% Authorship: Claude-drafted, author-revised
\section{Conclusion}

Where do the determinants of legal content come from? On Berman's answer
they come from practice, and this Article has taken that answer literally
enough to build it, run it, and measure what it delivers. Two things
followed, and they pull in different directions.

First, the machinery works. Given faithful weights, the
weighted sum returns one verdict where a rule reading each authority
only for validity returns two. On directly opposed cases the rule contradicts itself every time and the sum
never does, and the sum's half of that is an entailment of the aggregation
form rather than a lucky run. It absorbs weight error of about $\pm 0.09$, and better than twice that
under the correlated activations of realistic disputes, where a paired run
rises from $0.069$ to $0.156$. It reads how strongly a principle applies
rather than where the principle ranks. And against a tiered Hartian
hierarchy it leaves less of the law unsettled, by $+0.272$, positive on
every seed.

That last comparison carries a condition that belongs here
rather than only in Part IV, because it is the one most easily mislaid. The
margin is measured on independent activations. Under the correlated
activations this Article treats as realistic it reverses: the tiered
hierarchy decides $95.6\%$ of cases against the weighted sum's $90.1\%$.
The reversal does not show that a tiered rule finds more law. It shows that
such a rule decides faster by consulting less, bringing $1.83$ of nine
principles to bear and setting aside roughly three-quarters of the force in
play. Determinacy counted on its own rewards foreclosure. That is a reason
to ask what a rule hears before asking how much it settles.

Second, none of this fixes what the law is. Between them, four intelligible
readings of officials' taking-up of a principle produce four weight
structures, three of them distinct. Run over a common battery,
those structures return different verdicts on $36.1\%$ of cases on the bare
sign of net force, on $44.5\%$ once the determinacy threshold is applied, and
opposite determinate verdicts on $27.2\%$. Each clears every requirement the theory imposes, three of them on every
seed and appointment on twenty-six of thirty, its noise tolerance
straddling the floor. They differ not in competence but in
object: uptake reads what officials say, reliance and appointment read
what their decisions leave standing.

Put together, the two findings say something a quantity of law cannot
capture. On this evidence principled positivism is more determinate than Hart and
less settled than it looks. Its determinacy is determinacy under a construal,
and the theory does not fix the construal whose determinacy it is. The mark of a verdict its rivals would reverse is present in the output of
any single reading but unreliable, and it is worth least under precisely
the construal that discards the moral principles. Hart's gaps announce themselves to the judge standing in one.
This contingency does not.

Beyond that, several things one might read in are not in it. No verdict
reported here has been shown to be the wrong verdict, since the design
contains no standard that any of them could fail. For Berman, officials' taking-up of a principle makes its weight what it
is. That grounding relation is intelligible under every mechanism, and
under every mechanism it works. The prior question remains open: what
drives those takings-up over time. That is a gap in the
dynamics, not a flaw in the relation.

On the theory's side three positions remain available, and this Article has
priced them rather than choosing among them. First, a defender may single out one dynamic and defend it as the one real
systems run on. That converts a conceptual claim into an empirical bet.
Second, a defender may deny that the dynamics come apart where it matters.
That denial requires showing a convergence rather than assuming one. Or a defender may accept the indexing and its cost,
a relativity about legal content that positivists have generally wanted to
avoid.

Each is a substantive commitment about how legal systems in fact work, and
a commitment of that kind is the theory's to make rather than a
simulation's. A model of this sort can say what each option
costs. It cannot tell anyone which dynamic actually drives the
takings-up in a real legal system, because the practice it reads is the
practice it built.

So the work that follows is empirical, and it has a shape. The first task
is coding, assigning activation values to the principles at work in real
judicial opinions, so that the aggregation can be run against decided cases
rather than synthetic ones. Berman has already coded one case this way,
\textit{TVA v.\ Hill}, showing both that the exercise is possible and
how much judgment it takes. Whether the exercise scales past a handful of
opinions is the question that would have to be settled first.

The second task is discrimination among the dynamics, and it is tractable
because they leave different traces. Against the two mechanisms
that read holdings, uptake stands at rank correlations of $-0.78$ and
$-0.75$, where those two sit at $+0.93$ with each other. Real doctrine either exhibits that pattern
or it does not. A body of case law where a principle's authority tracks
how often it is invoked would look measurably unlike one where
authority tracks the holdings still standing, and the difference should be
visible to anyone willing to count.

A third implication runs the other way, from the model back toward
institutions. The emergent weights of Part V trace to the composition of the
bench rather than to any winnowing of judges as such. Change the mix of
methodologies enough and the line between the principles that carry weight and
those that carry none moves with it: give the methodology that reaches for the
moral principles a majority of the seats rather than a quarter, and three of
them cross the line while three textual and structural principles cross back. On this account, who sits on a court
would shape not only how particular cases come out but which principles
count as legally weighty at all. That is a stronger claim about judicial
selection than the usual one, and a testable one.

The limits travel with all of it, and there are more than two. The cases
are synthetic and the practice is simulated, so which construal describes real
adjudication is a question no number here answers. Four
mechanisms are a sample of the readings a practice-based theory could take
of its own key term rather than the space of them; a larger sample could
add structures to the three found here, though it could not remove them,
since each was built and run.

Three further limits are specific rather than general, and each attaches to
a particular claim. The margin over the hierarchical Hart holds on independent
activations
and reverses on correlated ones, so it is a claim about a case distribution
and not about the rules alone. The separation between weighing and counting
holds across every weight regime and every activation structure short of
full correlation, where a magnitude probe has almost nothing left to move.
And the tier sizes contribute to the divergence: balancing them at two
principles apiece takes the four-way rate from $0.361$ to $0.303$, and
cutting the moral tier to two takes it to $0.249$, so roughly a sixth to a
third of the divergence is attributable to composition, with the remainder
unexplained by it. No run reported here separates that remainder from the
category-size effect operating through correlated activations.

What remains, then, is narrower than the problem principled positivism
inherited. Between them, Hart left the law running out in hard cases and Dworkin left
the right answer resting on moral interpretation, and Berman's aggregation
genuinely improves on both. The question it puts in their place is which
facts about a single practice do the grounding, with a measurable
consequence attached to every answer. That is the problem of practice. As problems go it is a better one than it replaces, and it is now specified
closely enough to be worked on.
